\scanpage{22}
``définie sur $K_i$ pour $i$ assez grand''.
Choisissons un tel $K=K_i$, et essayons d'exprimer l'information supplémentaire
de type ``galoisien'' \unsure{qui l'exprime|qu'il exprime}. Soit $\overline K$ la clôture alg.
de $K$ dans $k$. Considérons
\[
(7)\quad\diagram{\begin{tikzcd}[column sep=large]
\mathcal X_{g,\nu\,\mathbb Q}\arrow[d]&\mathcal X_{g,\nu\,K}\arrow[l]\arrow[d,no head]&\mathcal X_{g,\nu\,\overline K}\arrow[l,no head]\arrow[d,no head]&\mathcal X_{g,\nu\,k}\arrow[l,no head]\arrow[d,no head]\\
M_{g,\nu\,\mathbb Q}\arrow[d]&M_{g,\nu\,K}\arrow[l]\arrow[d,no head]&M_{g,\nu\,\overline K}\arrow[l,no head]\arrow[d,no head]&M_{g,\nu\,k}\arrow[l,no head]\arrow[d,no head]\\
\mathbb Q&K\arrow[l,no head]\arrow[u,dashed,bend right=35]\arrow[uu,dashed,bend right=55]&\overline K\arrow[l,no head]&k\arrow[l,no head,"{\underbrace{\hspace{2.3cm}}_{\substack{\text{du pt de vue}\\\text{homotopique, cette}\\\text{extension est de}\\\text{nature triviale}}}}"']
\end{tikzcd}}
\]
$(X,x)$ provient d'une \deleted{\unsure{pt}}\inserted{section} (notée encore $x$)
de $\unsure{\mathcal X_{g,\nu\,K}|M_{g,\nu\,K}}$ sur $K$, définissant une section
(notée encore $\xi$) de $M_{g,\nu\,K}$ sur $K$, (dont sont déduites les sections
$x,\xi$ sur $k$, par changement de corps de base).

Posons
\[
(\unsure{8})\quad\left\{\begin{aligned}
\mathcal M_{X,x,K}&=\pi_1(\unsure{\mathcal X_{g,\nu\,K}|M_{g,\nu\,K}},\deleted{\xi}\inserted{x})\\
\mathcal N_{X,x,K}&=\pi_1(M_{g,\nu\,K},\xi)\\
\Gamma_{\overline K/K}&=\operatorname{Gal}(\overline K/K)=\pi_1(\operatorname{Spec}K,\unsure{\overline K|\overline k})
\end{aligned}\right.
\]
on a cette fois-ci, comme (1) avec $K$ remplaçant $\mathbb Q$ :
\mspage{21}
